\section{Ordinary tiles and the popular forest}\label{tile:section}

The purpose of this section is to separate the part of the packet
form controlled by ordinary size and density estimates from the part
requiring phase averaging. The size--density--tree decomposition follows the time-frequency
framework of Lacey and Thiele for the Carleson operator
\cite{LaceyThiele2000} and its planar directional adaptation by
Lacey and Li \cite{LLAnnular,LLStein}. We prove the
estimates in the present normalization, including arbitrary subcollections
and the Hilbert-valued summation needed for narrow vertical bins.

For the ordinary estimates, write the packet form \eqref{eq:src-3} as
\(\Lambda_{\mathcal S}(f,g)\), where \(\mathcal S\) is a finite
collection from one of the forms constructed in the preceding section.
The constants in these estimates depend on fixed packet adaptation,
separation, and multiplicity constants only.  In particular, they do
not depend on the number of tiles or lengths, on \(w\), or on
\(m=W/w\geq1\).  These ordinary estimates hold for every measurable
real-valued slope function \(u\).

\subsection{Density and the model estimates}

Fix one of the angular dyadic grids.  Denote by \(\Omega(t,l)\)
the interval in this grid at the level corresponding to the dyadic
length \(l\) that contains \(t\).  The intervals are half open.  Their
length is a fixed constant times \(w/l\); in particular,
\(\Omega(t,\ell_s)=\Omega_s\) whenever \(t\in\Omega_s\).

For a box \(R\) of length \(l\), width \(W\), slope \(t\), and center
\(z_R=(x_R,y_R)\), put
\begin{equation}\label{tile:eq-density-weight}
 B_R(x,y)=
 \left(1+\frac{|x-x_R|}{l}
 +\frac{|(y-y_R)-t(x-x_R)|}{W}\right)^{-D_0}.
\end{equation}
Here \(D_0\) is a sufficiently large fixed integer, chosen once below.
For a measurable set \(E\) of finite measure define
\begin{equation}\label{tile:eq-density}
 D_E(s)=
 \sup_{\substack{\ell_s\leq l\leq\tau,\ l\ \text{dyadic}\\
                   t\in\Omega_s,\ z_s\in R}}
 \frac{1}{|R|}
 \int_{E\cap\{u\in\Omega(t,l)\}} B_R(z)\,\dd z.
\end{equation}
The boxes in this supremum have the dimensions and slope used in
\eqref{tile:eq-density-weight}.  Since \(\int B_R\leq C|R|\), all
densities are bounded by a fixed constant.  Density is monotone in
\(E\); passing to a subcollection of tiles does not change the density
of any remaining tile.

\begin{proposition}[Ordinary model estimates]\label{tile:prop-model}
Let the bin Hilbert space be \(\mathcal H=\C^J\), where
\(J\leq C_{\rm bin}m\).  If \(\norm{g(z)}_{\mathcal H}\leq\one_E(z)\)
and \(D_E(s)\leq\gamma\) for every \(s\in\mathcal S\), with
\(0<\gamma\leq C\), then
\begin{equation}\label{eq:src-4}
 |\Lambda_{\mathcal S}(f,g)|
 \leq C(1+\log(2m))^C\gamma^{1/4}
             |E|^{1/2}\norm{f}_{L^2(\mathcal H)}.
\end{equation}
In the scalar case, if \(|f|\leq\one_F\), \(|g|\leq\one_E\), and
\(0<|E|<|F|<\infty\), then
\begin{equation}\label{tile:eq-reverse-disparity}
 |\Lambda_{\mathcal S}(f,g)|
 \leq C\sqrt{|F||E|}
                \left(\frac{|E|}{|F|}\right)^{9/20}.
\end{equation}
Both conclusions hold uniformly for every subcollection of
\(\mathcal S\).
\end{proposition}

We first prove the scalar estimates, suppressing the bin index \(j\);
size and density selection and the tree bound provide the scalar model
bounds. The final component summation gives the stated logarithmic
vector loss in \(m\).

\subsection{Maximal functions and separated synthesis}

For a slope \(t\), let \(\mathcal M_t\) be the uncentered strong
maximal operator in the coordinates \((x,y-tx)\).  We shall use
\begin{equation}\label{tile:eq-strong-maximal}
 \norm{\mathcal M_t h}_2\leq C\norm{h}_2,
\end{equation}
with a constant independent of \(t\).  Here and below maximal
operators act on the absolute value of their argument.  To recall a
proof, the one-dimensional uncentered maximal operator \(M\) is of
weak type \((1,1)\): a finite interval covering of a compact subset
of \(\{Mh>\lambda\}\) has a disjoint subfamily whose threefold
enlargements cover that subset, giving measure at most
\(3\lambda^{-1}\norm{h}_1\).  Approximation gives the full
inequality.  Apply it to
\(h\one_{\{|h|>\lambda/2\}}\), since the complementary part has
maximal function at most \(\lambda/2\), and integrate the
distribution estimate against \(2\lambda\,\dd\lambda\).
Fubini's theorem gives \(\norm{Mh}_2^2\leq C\norm{h}_2^2\).
The strong maximal operator is bounded pointwise by the composition
of the two one-dimensional maximal operators.  Fubini and the
determinant-one shear prove \eqref{tile:eq-strong-maximal}.

A consequence we use repeatedly is smooth reproduction on a box.
Suppose a Schwartz function, after removal of a fixed carrier, has
Fourier support
\[
 |\xi+t\eta|\leq A/r,\qquad |\eta|\leq A/W.
\]
A product of smooth low-pass cutoffs that equals one on this set
reproduces the function.  Its kernel is bounded, for any fixed
\(N\), by
\[
 \frac{C_{A,N}}{rW}
 \left(1+\frac{|x|}{r}\right)^{-N}
 \left(1+\frac{|y-tx|}{W}\right)^{-N}.
\]
Decomposing this majorant into dilated rectangles proves the
following statement.  If \(Q\) is an \(r\times W\) box of slope
\(t\), and \(F_\alpha\) is any finite family with these common
spectral bounds, then
\begin{equation}\label{tile:eq-reproduction}
 \sup_{z\in Q}\sup_\alpha |F_\alpha(z)|
 \leq C_A\inf_{z\in Q}
             \mathcal M_t\bigl(\sup_\alpha|F_\alpha|\bigr)(z).
\end{equation}
Indeed, replacing the center of any kernel rectangle by any other
point of \(Q\) changes its required enlargement by only an absolute
factor.  A common carrier has modulus one and does not affect this
argument.

A test tree has parameters \(T=(z_T,t,l)\), where \(l\leq\tau\)
is dyadic, and a box \(R_T\) of length \(l\), width \(W\), and
slope \(t\).  Its tiles satisfy
\begin{equation}\label{tile:eq-tree}
 \ell_s\leq l,\qquad z_s\in R_T,\qquad t\in\Omega_s.
\end{equation}
An arbitrary subset satisfying these conditions is also a tree.
Choose a fixed \(c_1>0\) so that
\(c_*\ll c_1\ll\min(c_b,c_\Omega)\), where \(c_b d_s\) is the lower distance
of the selector support from \(\theta_s\).  A tile is separated
in this tree if \(|t-\theta_s|\geq c_1d_s\).

For a positive integer \(J_0\) define the top polynomial
\begin{equation}\label{tile:eq-top-polynomial}
 \mathfrak p_{T,J_0}(z)=
 \left(1+\frac{(x-x_T)^2}{l^2}
 +\frac{((y-y_T)-t(x-x_T))^2}{W^2}\right)^{J_0}.
\end{equation}

\begin{lemma}[Separated synthesis]\label{tile:lem-synthesis}
Let \(\mathcal T\) be separated tiles satisfying
\eqref{tile:eq-tree}, and let \(\Upsilon_s\) be either of their
packet families \(\Phi_s\) or \(\Psi_s\).  For every fixed
\(J_0\),
\begin{align}
 \left\|\mathfrak p_{T,J_0}
                \sum_{s\in\mathcal T}c_s\Upsilon_s\right\|_2
 &\leq C_{J_0}
          \left(\sum_{s\in\mathcal T}|c_s|^2|R_s|\right)^{1/2},
                  \label{tile:eq-synthesis}\\
 \left\|\mathfrak p_{T,J_0}
       \sup_{a<b}\left|
        \sum_{\substack{s\in\mathcal T\\a<\ell_s\leq b}}
                  c_s\Upsilon_s\right|\right\|_2
 &\leq C_{J_0}
          \left(\sum_{s\in\mathcal T}|c_s|^2|R_s|\right)^{1/2}.
                  \label{tile:eq-maximal-synthesis}
\end{align}
The same statements hold for Hilbert-valued sums
\(\bigl(\sum_s c_{s,j}\Upsilon_{s,j}\bigr)_j\), with
\(\sum_s\norm{c_s}_{\mathcal H}^2|R_s|\) on the right and
constants independent of \(J\).  Bounded scalar factors depending
on \(s,j\), such as \(b_{s,j}(t)\), can be incorporated into
the coefficients.
\end{lemma}

\begin{proof}
From the packet support and \(t\in\Omega_s\),
\[
 \xi+t\eta=(\xi+\theta_s\eta)+(t-\theta_s)\eta,
 \qquad
 \frac{a_0}{\ell_s}\leq|\xi+t\eta|
                         \leq\frac{a_1}{\ell_s},
\]
where \(0<a_0<a_1<\infty\) are fixed.  Positivity of \(a_0\)
uses \(c_*\ll c_1\) and the two fixed bounds on \(w\eta\).
For envelopes, replacing slope \(\theta_s\) by \(t\) changes
the transverse coordinate over one tile length by at most
\(C d_s\ell_s=Cw\leq CW\).  Thus every fixed Schwartz
envelope remains a uniform envelope in the \(t\)-coordinates.

At one length there are boundedly many labels, since all labels
of the tree lie in an interval of length \(O(d_s)\).  Within
each label the centers lie on a translated box lattice.  The
inner product of two normalized packets at this length is
bounded by a fixed rapidly decreasing function of their
normalized center separation.  Its row and column sums are
bounded.  The elementary Schur inequality, obtained by applying
\(2|a_sa_q|\leq |a_s|^2+|a_q|^2\) to these row and column
sums, proves the fixed-length synthesis bound.

This argument also applies after multiplying packets by
\(\mathfrak p_{T,J_0}\).  Their centers belong to \(R_T\),
their lengths are at most \(l\), and their widths are \(W\).
Consequently the polynomial can be absorbed into a fixed
larger number of packet decay powers.  Multiplication by a
polynomial differentiates the Fourier transform and does not
enlarge its support.

Choose a fixed integer \(q_0\) such that
\(2^{q_0}a_0>4a_1\), and split dyadic lengths according to
their exponents modulo \(q_0\).  Within each resulting list,
the displayed annuli are disjoint and have gaps.  Orthogonality
in the first \(t\)-coordinate therefore proves
\eqref{tile:eq-synthesis}, also for the polynomially weighted
packets.

On one such list, the sum over all lengths above a given
endpoint is exactly a smooth low-pass convolution in that
coordinate applied to the complete sum: choose the cutoff to
equal one on the lower-frequency annuli retained at that
endpoint and zero on all higher-frequency annuli.  The gap
permits dilates of one fixed smooth cutoff.  Their absolute
kernels are dominated by the one-dimensional maximal
operator.  An interval of lengths is a difference of two
endpoint sums.  Apply this observation to the full weighted
sum, use its just-proved \(L^2\) bound, and sum over the fixed
number of lists.  This gives
\eqref{tile:eq-maximal-synthesis}.

For Hilbert-valued sums, the fixed-length and orthogonality
bounds follow by summing the scalar squared bounds over
components.  The same low-pass convolution acts on every
component, since \(a_0,a_1\) are common.  Its norm is
pointwise bounded by the positive kernel acting on the
Hilbert norm of the full sum.  The maximal-function proof
is therefore dimension free.  Alternatively one can sum
the squared scalar maximal estimates.
\end{proof}

For scalar coefficients
\(v_s=\langle f,\Phi_s\rangle/|R_s|\), define their size on
a collection by
\begin{equation}\label{tile:eq-size}
 \operatorname{size}(v)=
 \sup_T
 \left(\frac{1}{|R_T|}
       \sum_{\substack{s\in T\\|t-\theta_s|\geq c_1d_s}}
                       |v_s|^2|R_s|\right)^{1/2}.
\end{equation}
A one-tile test gives \(|v_s|\leq C\operatorname{size}(v)\).
If \(\Phi_s=0\), then \(v_s=0\). Otherwise choose
\((\xi,\eta)\in\supp\widehat\Phi_s\) and put \(q=-\xi/\eta\).
The angular margin and \(c_1\ll c_\Omega\) place both
\(q-2c_1d_s\) and \(q+2c_1d_s\) in \(\Omega_s\).
At least one is at distance at least \(2c_1d_s\) from \(\theta_s\);
use it as the direction of a top centered at \(z_s\) with length
\(\ell_s\). This is a separated one-tile test of area \(|R_s|\),
which proves the bound.
Moreover
\begin{equation}\label{tile:eq-bounded-size}
 \operatorname{size}(v)\leq C\norm{f}_\infty.
\end{equation}
To prove this, for the separated portion of a fixed tree put
\(B_T=\sum_s|v_s|^2|R_s|\) and
\(A_T=\sum_s v_s\Phi_s\).  The pairing is linear in its first
argument, so \(B_T=\langle f,A_T\rangle\).  By
\Cref{tile:lem-synthesis} and Cauchy--Schwarz with
\(\mathfrak p_{T,1}^{-1}\),
\[
 B_T\leq\norm{f}_\infty\norm{A_T}_1
 \leq C\norm{f}_\infty |R_T|^{1/2}B_T^{1/2}.
\]
Division, with the zero case understood, proves
\eqref{tile:eq-bounded-size}.

\subsection{Removing large size}

\begin{lemma}[Size selection]\label{tile:lem-size}
Suppose \(\operatorname{size}(v)\leq C_*\sigma\), with
\(C_*\) fixed.  One can remove a collection of trees so that
the remaining size is at most \(\sigma\) and
\begin{equation}\label{tile:eq-size-budget}
 \sum_T|R_T|\leq C\sigma^{-2}\norm{f}_2^2.
\end{equation}
The removed tiles partition into ordinary trees satisfying
\eqref{tile:eq-tree}; their total top area has the same bound.
\end{lemma}

\begin{proof}
First consider positive separated tests, with
\(t-\theta_s\geq c_1d_s\).  While a test has energy larger
than \(\sigma^2|R_T|/2\), choose one whose direction is within
\(\varepsilon\) of the infimum of the directions of all such
tests.  Here \(\varepsilon>0\) is a sufficiently small fixed
multiple of \(\min_{s\in\mathcal S}d_s\).  Delete all remaining
tiles with \(\ell_s\leq l_T\), \(t_T\in\Omega_s\), and
\(z_s\in A_0R_T\), including nonseparated tiles.  The fixed
absolute dilation \(A_0\) is chosen below.  Record the
positive separated test tiles present at selection as
\(\mathcal Q_T\).  The procedure terminates because every
selection deletes at least one tile.  Thereafter perform the
negative separated selection, using directions within
\(\varepsilon\) of the supremum and reversing the order.
At the end, each signed energy is at most
\(\sigma^2|R_T|/2\); hence the size is at most \(\sigma\).

We prove the area bound for the positive selection.  The
\(\mathcal Q_T\) are disjoint and
\begin{equation}\label{tile:eq-selected-energy}
 B:=\sum_T\sum_{s\in\mathcal Q_T}|v_s|^2|R_s|
       \asymp \sigma^2\sum_T|R_T|.
\end{equation}
The lower comparison is the selection threshold and the
upper one follows from the initial size bound, which
persists on every remaining subcollection.  We claim
\begin{equation}\label{tile:eq-selected-bessel}
 \left\|\sum_T\sum_{s\in\mathcal Q_T}v_s\Phi_s\right\|_2^2
 \leq CB.
\end{equation}

Expand the squared norm.  Pairs with lengths in a fixed
ratio have total contribution \(O(B)\).  Indeed, spectral
overlap forces their labels to lie within \(O(d_s)\), so
there are boundedly many neighboring labels.  Their box
shapes are comparable because the transverse discrepancy
over either length is \(O(w)\leq O(W)\).  Packet decay and
the spatial lattices then give the bounded row and column
sums for the normalized Gram matrix, exactly as in the
fixed-length proof above.

For the other pairs fix the top \(T\) of the longer member
\(q\in\mathcal Q_T\), and write \(t=t_T\).  Let
\(\mathcal A_T\) consist of the selected tiles \(s\) for which
\(\ell_q>C_0\ell_s\) and
\(\langle\Phi_s,\Phi_q\rangle\ne0\) for at least one such
\(q\).  Choose the fixed \(C_0\) sufficiently large.
Spectral overlap and the common vertical frequency range
give
\[
 |\theta_s-\theta_q|\leq Cc_*(d_s+d_q),\qquad
 |\theta_q-t|\leq C_\Omega d_q.
\]
Thus, by the choices of \(c_*\) and \(C_0\),
\begin{equation}\label{tile:eq-short-direction}
 |\theta_s-t|\leq \tfrac14 c_1d_s,
 \qquad t\in\Omega_s.
\end{equation}
The second assertion uses the fixed margin about the wave
slope support inside \(\Omega_s\).

We next verify a bounded overlap fact for the boxes of
\(\mathcal A_T\) in the coordinates of slope \(t\).  These
boxes can all be replaced by fixed comparable
\(\ell_s\times W\) boxes in those coordinates.  Suppose
two of the latter overlap, with \(\ell_s\) much smaller
than \(\ell_{s'}\).  Let \(t_s,t_{s'}\) be the directions
of their respective selecting tests.  Because \(s\) was
in a positive separated test,
\[
 t_s\geq\theta_s+c_1d_s>t+\tfrac12c_1d_s,
 \qquad
 |t_{s'}-t|\leq C d_{s'}.
\]
After separating the lengths by another fixed factor,
\(t_{s'}<t+\tfrac14c_1d_s\).  If the test selecting
\(s'\) had not yet been selected when \(s\) was selected,
all its eventual test tiles would still have been
present at that time.  It would therefore already
have been failing, since deleting tiles only
decreases each fixed test energy.
The near-infimum rule, with
\(\varepsilon\ll c_1\min d_s\), therefore forces the
test of \(s'\) to have been selected first.

Its length is at least \(\ell_{s'}\).  Overlap of the
comparable \(t\)-boxes bounds the horizontal center
separation by \(C\ell_{s'}\).  In the coordinates of
\(t_{s'}\), the transverse center separation is
\(O(W+w)=O(W)\): use
\(|t_{s'}-t|\ell_{s'}=O(w)\), or equivalently compare
each original tile over its own length.  Also
\(t_{s'}\in\Omega_s\), by the margin in
\eqref{tile:eq-short-direction}.  For an absolute
\(A_0\) this places \(s\) in the enlarged deletion tree
of \(s'\), a contradiction.  At comparable lengths,
bounded label and lattice multiplicity give bounded
overlap directly.  Splitting lengths into a fixed number
of lists proves the asserted overlap bound for all of
\(\mathcal A_T\).  The same reasoning applies to any
fixed enlargement of these boxes by increasing \(A_0\)
once.

Let \(B_s^*\) be sufficiently rapidly decreasing uniform
tile envelopes in the \(t\)-coordinates, and put
\(D_T=\mathfrak p_{T,J_0}^{-1}\), with \(J_0\) a fixed
large integer.  Then
\begin{equation}\label{tile:eq-short-envelopes}
 \left\|D_T\sum_{s\in\mathcal A_T}B_s^*\right\|_2
       \leq C|R_T|^{1/2}.
\end{equation}
Here centers of \(\mathcal A_T\) need not belong to
\(R_T\).  To prove the estimate, test against
\(\norm{h}_2\leq1\).  Since \(\ell_s\leq l_T\), the
top weight can be passed across a tile envelope by
spending fixed decay powers.  Decomposition into dilated
tile boxes, followed by the overlap just proved, gives
\begin{align*}
 \sum_s\int |h|D_T B_s^*
 &\leq C\sum_s |R_s|
       \inf_{R_s^t}D_T\inf_{R_s^t}\mathcal M_t h\\
 &\leq C\int D_T\mathcal M_t h
 \leq C\norm{D_T}_2\norm{h}_2
 \leq C|R_T|^{1/2}.
\end{align*}
The notation \(R_s^t\) denotes the fixed comparable
\(t\)-box; changing it by an absolute dilation is
harmless.  This proves \eqref{tile:eq-short-envelopes}.

For each \(s\in\mathcal A_T\), the sum of its cross
terms with \(\mathcal Q_T\) can include all
\(q\in\mathcal Q_T\) with \(\ell_q>C_0\ell_s\), since
the added inner products are zero.  Therefore its
absolute value is bounded using
\[
 U_T(z)=\sup_{r>0}
       \left|\sum_{\substack{q\in\mathcal Q_T\\\ell_q>r}}
                    v_q\Phi_q(z)\right|.
\]
The coefficient bound \(|v_s|\leq C\sigma\),
\eqref{tile:eq-short-envelopes}, and
\Cref{tile:lem-synthesis} show that all such cross terms
for this top have sum at most
\[
 C\sigma\int\left(\sum_{s\in\mathcal A_T}B_s^*\right)U_T
 \leq C\sigma |R_T|^{1/2}\norm{D_T^{-1}U_T}_2
 \leq C\sigma^2|R_T|.
\]
Summing and using \eqref{tile:eq-selected-energy}
proves \eqref{tile:eq-selected-bessel}.  Pairing its
synthesis with \(f\) gives
\(B\leq\norm{f}_2(CB)^{1/2}\), hence
\(B\leq C\norm{f}_2^2\).  The negative selection has
the identical proof with order reversed.

Finally, partition each enlarged deletion box
\(A_0R_T\) into a fixed number of translated ordinary
top boxes of length \(l_T\) and width \(W\), assigning
tiles by their centers.  Their direction and length
conditions are unchanged, and their total area is
\(O(|R_T|)\).  This proves the last assertion and
\eqref{tile:eq-size-budget}.
\end{proof}

\subsection{Removing large density}

\begin{lemma}[Density selection]\label{tile:lem-density}
For \(0<\eta\leq C\), all tiles with \(D_E(s)>\eta\)
can be assigned to keeper boxes \(R_T\), of dyadic
length \(l_T\leq\tau\), width \(W\), and slope \(t_T\),
with the following properties.  There are pairwise
disjoint measurable sets \(E_T\) such that
\begin{gather}
 E_T\subset E\cap K_\eta R_T
                 \cap\{u\in\Omega(t_T,l_T)\},\qquad
 |E_T|\geq c\eta|R_T|,
                  \label{tile:eq-density-witnesses}\\
 \sum_T|R_T|\leq C\eta^{-1}|E|,\qquad
 K_\eta\leq C(1+\eta^{-1/100}),
                  \label{tile:eq-keeper-budget}\\
 \ell_s\leq l_T,\qquad t_T\in\Omega_s,\qquad
 z_s\in C K_\eta R_T
 \quad\text{for each tile assigned to }T.
                  \label{tile:eq-density-assignment}
\end{gather}
After partitioning by centers, the assigned tiles form
ordinary test trees of total top area at most
\begin{equation}\label{tile:eq-density-budget}
 C\eta^{-1}(1+\eta^{-1/20})|E|.
\end{equation}
\end{lemma}

\begin{proof}
For each tile to be removed select one box in
\eqref{tile:eq-density} witnessing an average greater
than \(\eta\).  Since
\[
 \int_{\R^2\setminus KR}B_R\leq C K^{2-D_0}|R|,
\]
choosing the fixed \(D_0\) sufficiently large and then
\(K=K_\eta\leq C(1+\eta^{-1/100})\) retains at least
\(\eta|R|/2\) of the weighted test inside \(KR\).
Its unweighted tested set has measure at least
\(\eta|R|/2\), because \(B_R\leq1\).

Process these finitely many witnesses in decreasing
length order.  Keep a witness unless both its
\(K\)-dilated spatial box and its dyadic direction
interval intersect the corresponding objects of an
earlier keeper.  Assign it to any such keeper when
it is rejected.  The tested sets of the keepers are
disjoint: a common tested point would force both
intersections and preclude keeping the later one.
This proves \eqref{tile:eq-density-witnesses} and
\eqref{tile:eq-keeper-budget}.

Let a rejected witness have length \(l'\), direction
\(t'\), and original associated tile \(s\), and let
its keeper have length \(l_T\geq l'\) and direction
\(t_T\).  Dyadic nesting gives
\[
 \Omega(t_T,l_T)\subset\Omega(t',l')\subset\Omega_s,
 \qquad |t_T-t'|\leq Cw/l'.
\]
Choose a point common to the two dilated spatial
boxes.  Moving from it to the rejected box center,
or to \(z_s\), costs horizontally \(O(Kl')\) and
transversely in the \(t_T\)-coordinates
\(O(KW+K l'|t_T-t'|)=O(KW)\).  This proves
\eqref{tile:eq-density-assignment}.  A keeper's own
associated tiles satisfy the same conclusion.

Partition \(CKR_T\) into \(O(K^2)\) translated
boxes of length \(l_T\), width \(W\), and slope
\(t_T\), assigning tiles by their centers.
The length and angular requirements for a test
tree persist.  Their total area is at most
\(C(1+\eta^{-1/50})\eta^{-1}|E|\), which is bounded
by \eqref{tile:eq-density-budget} after changing
the absolute constant.
\end{proof}

\subsection{The tree estimate}

\begin{lemma}[Tree bound]\label{tile:lem-tree}
Suppose a scalar tree \(\mathcal T\) satisfies
\(D_E(s)\leq\eta\) on its tiles and has size at
most \(\sigma\).  If \(|g|\leq\one_E\), then
\begin{equation}\label{tile:eq-tree-bound}
 |\Lambda_{\mathcal T}(f,g)|
           \leq C\sigma\eta|R_T|.
\end{equation}
The size assumption can be replaced by the same
bound on any encompassing collection.
\end{lemma}

\begin{proof}
The assertion is immediate for an empty tree, so assume
the tree is nonempty. Work in centered coordinates \((x,y-tx)\) of the
top, so its coordinate extents are \(l\times W\).
For clarity write the transverse coordinate as
\(y'\), with top center \((0,0)\).  All tile
centers satisfy \(|x_s|\leq l/2\) and
\(|y_s'|\leq W/2\).  The packet envelopes are
uniform in these coordinates.

Partition the horizontal line into maximal dyadic
intervals \(I\) for which there is no tile with
\(x_s\in3I\) and \(\ell_s\leq |I|\).  The tiles
are finite and have positive minimum length, so
these intervals cover the line up to endpoints
and are pairwise disjoint.  Each parent fails
the condition, furnishing a tile \(q_I\) with
\begin{equation}\label{tile:eq-whitney-witness}
 \ell_{q_I}\leq2|I|,\qquad
 \dist(x_{q_I},I)\leq C|I|.
\end{equation}
Fix a sufficiently large absolute dyadic constant
\(C_{\mathrm W}\).  We treat lengths
\(\ell_s\leq C_{\mathrm W}|I|\) and longer lengths
separately on the strip above \(I\).

For a short tile, density tested at its own box
and wave decay give, for every fixed large \(N_0\),
\begin{equation}\label{tile:eq-short-tree-term}
 \int_{E\cap\{x\in I\}}
     |v_s b_s(u)\Psi_s|
 \leq C\sigma\eta|R_s|
        \left(1+\frac{\dist(x_s,I)}{\ell_s}\right)^{-N_0}.
\end{equation}
Indeed the selector support lies in \(\Omega_s\),
the density definition includes \(R_s\) with
direction \(\theta_s\), and a fixed number of
additional wave-decay powers extracts the last
factor while retaining the density weight.

Here is an explicit summation of
\eqref{tile:eq-short-tree-term}.  At a fixed
length there are \(O(1)\) labels, and each
horizontal column contains \(O(1)\) centers
within the transverse top strip.  If
\(\ell_s\leq r:=|I|\), the defining property
of \(I\) gives \(\dist(x_s,I)\geq r\).
Summing over the one-dimensional center lattice
then bounds its contribution, divided by
\(C\sigma\eta W\), by
\(C r(\ell_s/r)^{N_0-1}\).
For \(r<\ell_s\leq C_{\mathrm W}r\), the same
lattice sum is \(O(r)\) after multiplication
by \(\ell_s\), and there are only a fixed
number of such lengths.  Thus the total is
\(O(\sigma\eta Wr)\).
All intervals with \(r\leq l\) lie within a
fixed enlargement of the top's horizontal
interval, by \eqref{tile:eq-whitney-witness};
their disjoint lengths sum to \(O(l)\).
If \(r>l\), all tile lengths are at most
\(l\), all centers are at distance at least
\(r\) from \(I\), and their restriction
\(|x_s|\leq l/2\) improves the preceding
sum to
\(C\sigma\eta Wl(l/r)^{N_0-1}\).
For each such dyadic \(r\), only \(O(1)\)
intervals can occur, again by the parent
witness.  Summing proves a total short-tile
contribution \(C\sigma\eta lW\).

Now let \(I\) admit longer tiles.  Choose a
dyadic \(l'\asymp |I|\), greater than
\(\ell_{q_I}\) and small enough that
\(l'<\ell_s\) for every longer tile.
Increasing \(C_{\mathrm W}\) once ensures this.
It also ensures \(l'\leq\tau\).  By
\eqref{tile:eq-whitney-witness}, a box of
slope \(t\), length \(l'\), and width \(W\)
can contain \(z_{q_I}\) and the portion
\(I\times[-W/2,W/2]\), with a fixed
enlargement absorbed into the choice of
\(l'\).  Dyadic nesting puts every longer
selector support inside \(\Omega(t,l')\).
Consequently the original longer-tile sum vanishes outside
\(A_I=\{u\in\Omega(t,l')\}\).  We perform all following
freezing and error decompositions after multiplying that
sum by \(\one_{A_I}\), and retain this indicator on both
the main term and the error.  Individually the frozen
terms need not vanish outside \(A_I\).
For the cells
\[
 Q_{I,k}=I\times[(k-\tfrac12)W,(k+\tfrac12)W],
       \qquad k\in\Z,
\]
the density test at \(q_I\) therefore yields
\begin{equation}\label{tile:eq-cell-density}
 \mu_{I,k}:=
 |Q_{I,k}\cap E\cap\{u\in\Omega(t,l')\}|
 \leq C\eta(1+|k|)^{D_1}|I|W,
\end{equation}
where the fixed \(D_1\) can be taken to be
\(D_0\).  The density weight on the cell
is bounded below by a constant times
\((1+|k|)^{-D_0}\).

Put \(a=|u-t|\) at the point of evaluation.
For nonseparated tiles, selector support
implies \(a\asymp d_s\), so only a fixed
number of dyadic lengths occur.  Their
pointwise sum is bounded by
\(C\sigma(1+|k|)^{-N_1}\) on \(Q_{I,k}\),
for any fixed \(N_1\).  For separated
tiles replace \(b_s(u)\) by
\(\one_{\{d_s\geq a\}}b_s(t)\).
If \(d_s\geq a\), the error is at most
\(Ca/d_s\), using the derivative bound
on \(b_s\).  If \(d_s<a\), the actual
selector vanishes unless \(d_s\geq a/C\),
since \(|t-\theta_s|\leq C_\Omega d_s\).
There are only a fixed number of these
remaining scales.  Summing the geometric
series in \(a/d_s\), and using lattice
overlap at each scale, bounds the entire
error by \(C\sigma(1+|k|)^{-N_1}\).
This argument also covers \(a=0\), when
the separated error is zero.

The resulting main sum, at each point,
is an interval sum of the separated
packets with coefficients \(v_sb_s(t)\):
its lengths obey
\(C_{\mathrm W}|I|<\ell_s\leq w/a\).
Let \(M_{I,k}\) be the supremum on
\(Q_{I,k}\) of the absolute values of
all such interval sums with arbitrary
upper endpoint.  We claim that, for
every fixed sufficiently large \(N_2\),
\begin{equation}\label{tile:eq-cell-suprema}
 \sum_{I,k} M_{I,k}^2(1+|k|)^{N_2}|I|W
                    \leq C\sigma^2|R_T|,
\end{equation}
where only the \(I\) admitting longer
tiles occur.  To see this, let \(U\)
be the maximal interval sum over all
separated lengths of the tree.  For
each fixed \(I\), every interval sum
under consideration has horizontal
\(t\)-frequency at most \(C/|I|\)
and vertical frequency, after the
common carrier is removed, at most
\(C/W\).  The same is true after
multiplication by any top polynomial.
Apply \eqref{tile:eq-reproduction}
to the weighted sums and dominate
their absolute values by
\(\mathfrak p_{T,J_0}U\).  Since the
cells are disjoint, the squared cell
suprema are bounded in sum by
\(C\norm{\mathcal M_t
        (\mathfrak p_{T,J_0}U)}_2^2\).
For all the present cells
\(|x|\leq Cl\), and
\(\mathfrak p_{T,J_0}\asymp
 (1+|k|)^{2J_0}\) there.  Taking
\(4J_0\geq N_2\),
\eqref{tile:eq-strong-maximal} and
\Cref{tile:lem-synthesis} prove
\eqref{tile:eq-cell-suprema}.

The intervals admitting longer tiles
have total length \(O(l)\), by their
parent witnesses and
\(|I|<l/C_{\mathrm W}\).  Choose
\(N_2>2D_1+2\).  Factoring the full
density bound before Cauchy--Schwarz,
\begin{align*}
 \sum_{I,k}M_{I,k}\mu_{I,k}
 &\leq C\eta
   \left(\sum_{I,k}M_{I,k}^2
          (1+|k|)^{N_2}|I|W\right)^{1/2}\\
 &\hspace{1.2cm}\cdot
   \left(\sum_{I,k}(1+|k|)^{2D_1-N_2}|I|W\right)^{1/2}
 \leq C\sigma\eta|R_T|.
\end{align*}
The previously bounded pointwise
errors have the same total bound
by \eqref{tile:eq-cell-density},
taking \(N_1>D_1+2\).  Adding the
short contribution proves
\eqref{tile:eq-tree-bound}.
\end{proof}

\subsection{Proof of the model estimates}

\begin{proof}[Proof of \Cref{tile:prop-model}]
We first prove the scalar part of
\eqref{eq:src-4}.  If \(\norm{f}_2=0\)
or \(|E|=0\), the form is zero; assume
otherwise.  Process thresholds
\(\eta=2^{-n}\) in decreasing order,
starting high enough that no density
or size condition fails.  At each
step first remove densities exceeding
\(\eta\) by \Cref{tile:lem-density},
then lower size to
\[
 \sigma_\eta=\sqrt\eta\,
                    \norm{f}_2/|E|^{1/2}
\]
by \Cref{tile:lem-size}.  At the
start of the step, the remaining
collection has density at most
\(2\eta\), as well as the original
bound \(\gamma\), and size at most
\(\sqrt2\sigma_\eta\).  These
bounds persist under all removals
within the step.

The density removals have total
ordinary top area at most
\(C(1+\eta^{-1/20})\eta^{-1}|E|\),
and the size removals at most
\(C\sigma_\eta^{-2}\norm{f}_2^2
 =C\eta^{-1}|E|\).
Applying \Cref{tile:lem-tree} to
every resulting ordinary tree and
summing bounds the form by
\begin{equation}\label{tile:eq-threshold-sum}
 C\norm{f}_2|E|^{1/2}
 \sum_{\eta\ \text{dyadic}}
       \min(\gamma,2\eta)\sqrt\eta\,
             (1+\eta^{-1/20})\eta^{-1}.
\end{equation}
There is no limiting remainder to
estimate for a finite collection.
Every positive-density tile is
eventually removed.  If \(D_E(s)=0\),
testing at \(R_s\) shows that
\(b_s(u)g=0\) almost everywhere,
since the weight is strictly positive;
such a term is zero.  Zero-size
terms are zero as well.

For \(0<\gamma\leq1\), the terms
with \(\eta\leq\gamma\) sum to
\(C\gamma^{9/20}\), and the terms
with \(\eta\geq\gamma\) do also:
the relevant powers are
\(\eta^{1/2},\eta^{9/20}\) below
the crossover and
\(\gamma\eta^{-1/2},
  \gamma\eta^{-11/20}\) above it.
For bounded larger \(\gamma\) the
sum is uniformly bounded.  This
proves the asserted scalar
\(\gamma^{1/4}\) estimate, with
room in its exponent.

For \eqref{tile:eq-reverse-disparity},
put \(r=|E|/|F|<1\).  In addition
to the threshold size estimate,
every tree has size at most a fixed
constant by \eqref{tile:eq-bounded-size}.
Its actual size at step \(\eta\)
is therefore at most
\(C\min(1,\sqrt{\eta/r})\).
No change in the removal algorithm
is necessary.  With the universal
density bound, the sum of the tree
contributions at that step is at most
\[
 C|E|\min(1,\sqrt{\eta/r})
              \min(1,\eta)\eta^{-1}
                    (1+\eta^{-1/20}).
\]
For \(r\leq\eta\leq1\), this sums
to \(C|E|r^{-1/20}\); the scales
\(\eta>1\) cost \(C|E|\).
For \(\eta\leq r\), the sum is at most
\[
 C|E|r^{-1/2}
          \sum_{\eta\leq r}
               (\eta^{1/2}+\eta^{9/20})
 \leq C|E|r^{-1/20}
 =C\sqrt{|F||E|}\,r^{9/20}.
\]
This proves the reverse-disparity
claim, including the density-area
inflation from the selection.

Finally consider \(\mathcal H=\C^J\).
For a scalar measurable \(h\)
supported in \(E\), write
\[
 \norm{h}_{2,1}^{\#}
       =\int_0^\infty
             |\{|h|>a\}|^{1/2}\,\dd a.
\]
Layer-cake decomposition, retaining
the complex phase of \(h\), extends
the scalar estimate to
\[
 |\Lambda_j(f_j,g_j)|
       \leq C\gamma^{1/4}\norm{f_j}_2
                                  \norm{g_j}_{2,1}^{\#}.
\]
Every sliced support lies in \(E\),
so its density still has upper bound
\(\gamma\).  As \(|g_j|\leq1\),
split the last integral at
\(a_0=(2+m)^{-2}\).  The portion
below \(a_0\) is at most
\(a_0|E|^{1/2}\); for the rest,
Cauchy--Schwarz gives
\[
 \int_{a_0}^1|\{|g_j|>a\}|^{1/2}\,\dd a
 \leq
 \left(\int_{a_0}^1 a|\{|g_j|>a\}|\,\dd a\right)^{1/2}
 \left(\int_{a_0}^1\frac{\dd a}{a}\right)^{1/2}
 \leq C\sqrt{1+\log(2m)}\norm{g_j}_2.
\]
Since \(J\leq C_{\rm bin}m\) and
\(\sum_j\norm{g_j}_2^2\leq|E|\),
\[
 \sum_j(\norm{g_j}_{2,1}^{\#})^2
       \leq C(1+\log(2m))|E|.
\]
Cauchy--Schwarz in \(j\) proves
\eqref{eq:src-4}, even with a square
root of this logarithmic factor.
All arguments used only the stated
packet hypotheses, which persist
on every subcollection.  Every
choice of decay or derivative order
above is fixed before the parameters
and the finite tile collection are
chosen.  This completes the proof
of \Cref{tile:prop-model}.
\end{proof}

For later parameter choices we record a fixed numerical loss
available from this proof.  In the regime
\(m\asymp\exp(L^{0.9})\), with \(L\) sufficiently large,
the last vector estimate implies
\begin{equation}\label{tile:eq-improvement-loss}
 |\Lambda_{\mathcal S}(f,g)|
 \leq C L^{b_4}\gamma^{c_4}|E|^{1/2}
                         \norm{f}_{L^2(\mathcal H)},
 \qquad b_4=1,\quad c_4=\tfrac14.
\end{equation}
These exponents are fixed independently of all subsequent density,
forest, and mask choices.
